Policy-aware data exchange systems increasingly need to move large or continuously produced datasets across independently governed services. Traditional request-response communication is simple to reason about, but it often forces complete result materialisation before downstream processing can begin. Streaming communication can reduce this bottleneck by transferring data incrementally, yet it also changes failure behaviour, buffering, backpressure, and the relationship between the data path and policy-controlled orchestration.

This thesis studies how streaming communication can be introduced into dynamically composed, policy-aware microservice systems without weakening their control model. DYNAMOS, a middleware for policy-driven data exchange developed at the University of Amsterdam, is used as the case-study system because it combines generated service chains, data-sharing archetypes, sidecar-mediated communication, and explicit policy approval. The research first compares candidate transport mechanisms in a controlled Kubernetes pipeline, using throughput, latency, pacing, correctness, recovery, and resource pressure as evaluation dimensions. It then implements three chunk-aware DYNAMOS data paths: chunked unary transfer, intra-job gRPC streaming, and RabbitMQ Streams at the agent boundary.

A local validation on bulk \texttt{dataThroughTtp} retrieval shows that all three chunk-aware implementations reduce completion time and peak memory compared with the original classic-unary baseline: at 250{,}000 returned rows, completion time drops by a factor of 2.1--2.5 (29.1\,s to 11.7--14.1\,s) and peak sampled memory drops from 3.7\,GiB to 1.4--1.7\,GiB, with additional checks on \texttt{computeToData} and multi-provider \texttt{dataThroughTtp}. The results show that the main improvement comes from preserving chunks through the data path rather than from streaming as a label. Direct gRPC transfer is strongest when service lifetimes are naturally coupled; brokered streams remain valuable when decoupling, retention, and replay matter. The broader conclusion is that streaming should be treated as a post-policy data-path change: the approved execution path remains policy-governed, while large result transfer becomes incremental and more scalable.
